% Mejora de saltos de página y tipografía (plantilla UNIR)
\let\maketitle\relax

% Imágenes: escalado para que los diagramas quepan en la página
\usepackage{graphicx}
\setkeys{Gin}{width=0.88\linewidth,keepaspectratio}

% Encabezado y pie de página (plantilla UNIR)
\usepackage{fancyhdr}
\fancyhf{}
\fancyhead[C]{\normalfont\small Rehabilitación motriz asistida por visión por computadora}
\fancyfoot[C]{\thepage}
\renewcommand{\headrulewidth}{0.4pt}
\pagestyle{fancy}
\thispagestyle{empty}

% Evitar viudas y huérfanas (líneas sueltas)
\widowpenalty=10000
\clubpenalty=10000

% Mantener títulos con el párrafo siguiente
\usepackage{titlesec}
\titleformat{\section}{\normalfont\Large\bfseries}{\thesection}{1em}{}
\titleformat{\subsection}{\normalfont\large\bfseries}{\thesubsection}{1em}{}
\titleformat{\subsubsection}{\normalfont\normalsize\bfseries}{\thesubsubsection}{1em}{}
\titlespacing*{\section}{0pt}{2.5ex plus 1ex minus .2ex}{1.5ex plus .2ex}
\titlespacing*{\subsection}{0pt}{2.25ex plus 1ex minus .2ex}{1.25ex plus .2ex}
\titlespacing*{\subsubsection}{0pt}{2ex plus 1ex minus .2ex}{1ex plus .2ex}

% Tablas y figuras: título arriba, fuente abajo (plantilla UNIR / rúbrica)
\usepackage{caption}
\captionsetup{
  tableposition=top,
  figureposition=bottom,
  labelfont=bf,
  font=small,
  skip=4pt
}
\captionsetup[table]{position=above}
\captionsetup[figure]{position=below}

% Párrafos: sin sangría, espacio entre párrafos (estilo plantilla)
\setlength{\parindent}{0pt}
\setlength{\parskip}{6pt}

% Mejor corte de líneas
\sloppy
